\chapter{绪论}
\setcounter{chapter}{1}
\chapterintro{
掌握环境卫生学的定义、研究对象和内容；了解我国环境卫生工作主要成就及全球环境问题；理解"绿水青山就是金山银山"的发展理念。
}

% ----- 1.1 概述 -----
\section{环境卫生学概述}

\begin{defbox}
\textbf{环境卫生学（Environmental Hygiene / Environmental Health）}是研究自然环境和生活环境与人群健康的关系，揭示环境因素对人群健康影响的发生、发展规律，为充分利用环境有益因素和控制环境有害因素提出卫生要求和预防对策，增进人体健康，维护和提高人体健康水平的学科。
\end{defbox}

\subsection{研究对象与内容}

环境卫生学以\keyword{人类及其周围的环境}为研究对象，包括自然环境与生活环境。

\begin{tcolorbox}[colback=light, colframe=main, boxrule=0.5pt, arc=4pt, title=\textbf{环境卫生学的研究内容}, breakable]
\begin{enumerate}[leftmargin=*, itemsep=3pt]
    \item \imp{环境与健康关系的基础理论研究}——揭示环境因素对健康的影响机制。包括：
    \begin{itemize}
        \item \textbf{人类基因组计划（HGP）}：完成了人类23对染色体约60亿个核苷酸排列顺序的测定
        \item \textbf{环境基因组计划（EGP）}：研究有重要功能的环境应答基因多态性，确定引起环境暴露致病危险性差异的遗传因素
        \item \textbf{国际人类基因组单体型图谱（HapMap）}：构建不同人群的高密度单核苷酸多态（SNP）图谱，揭示群体分子遗传机制
    \end{itemize}
    \item \imp{环境因素与健康关系的确认性研究}——运用流行病学和毒理学方法。环境因素对健康影响的特点：\textbf{长期、低剂量、反复作用}
    \item \imp{创建和引进环境卫生学研究的新技术和新方法}——提升研究精度与深度
    \item \imp{研究环境卫生监督体系的理论依据}——为卫生标准与法规提供科学支撑
\end{enumerate}
\end{tcolorbox}

\subsection{环境的基本概念}

\begin{defbox}
\textbf{环境（Environment）}：指作用于人类所有外界力量（因素）的总和。人与环境既相互对立，又相互制约，既相互依存又相互转化，两者存在对立统一关系。
\end{defbox}

\begin{defbox}
\textbf{环境介质（Environmental Media）}：人类赖以生存的物质环境条件，通常以气态、液态和固态三种物质形态存在，能够容纳和运载各种环境因素。主要指\keyword{大气、水、土壤（岩石）以及包括人体在内的所有生物体}。
\end{defbox}

\textbf{环境介质的特点：}
\begin{enumerate}[leftmargin=*]
    \item 三种物质形态往往不完全以单一介质形式存在
    \item 三种物质形态在一定条件下可以相互转化
    \item 环境介质的运动可携带环境污染物向远方扩散
    \item 环境介质还能维持自身的稳定状态
\end{enumerate}

\begin{defbox}
\textbf{环境因素（Environmental Factors）}：被介质容纳和转运的成分或介质中各种无机和有机的组成成分。分为\keyword{物理性、化学性、生物性}三类。

\begin{itemize}[leftmargin=*]
    \item \textbf{物理因素：}小气候（气温、气湿、气流和热辐射，对机体热平衡产生明显影响）、噪声、振动、非电离辐射、电离辐射
    \item \textbf{化学因素：}无机和有机化学物
    \item \textbf{生物因素：}细菌、真菌、病毒、寄生虫和生物性变应原（如植物花粉、真菌孢子等）
\end{itemize}
\end{defbox}

\subsection{原生环境与次生环境}

\begin{defbox}
\textbf{原生环境（Primary Environment）}：天然形成的未受或少受人为因素影响的环境。有利方面：清洁的空气和水、充足的日光、适宜的小气候、秀丽的自然风光。不利方面：自然灾害、天然植物毒素、天然动物毒素、\textbf{生物地球化学性疾病}。

\textbf{次生环境（Secondary Environment）}：受人为活动影响形成的环境。有利方面：水利工程、资源开采、化工行业发展等。不利方面：污染环境，影响健康——全球气候变暖、臭氧层破坏、酸雨、生物多样性锐减；\textbf{环境污染性疾病}。
\end{defbox}

\begin{defbox}
\textbf{环境内分泌干扰物（EDCs）}：许多环境化学污染物（如有机氯化合物、二噁英、烷基酚、邻苯二甲酸酯等）对维持机体内环境稳态和调节发育过程的体内天然激素的生成、释放、转运、代谢、结合、效应造成严重影响的外源性物质，被称为内分泌干扰化学物。
\end{defbox}

% ----- 1.2 污染物 -----
\section{污染物分类}

\begin{defbox}
\textbf{一次污染物（Primary Pollutant）}：由污染源直接排入环境未发生变化的污染物。
\end{defbox}

\begin{defbox}
\textbf{二次污染物（Secondary Pollutant）}：一次污染物进入环境后在物理、化学或生物学作用下，或与其他物质发生反应而形成与原来污染物的理化性质和毒性完全不同的新污染物。
\end{defbox}

\begin{keybox}
常见的二次污染物：光化学烟雾中的臭氧、过氧乙酰硝酸酯（PAN）、酸雨中的硫酸和硝酸等。二次污染物的毒性往往比一次污染物更大。
\end{keybox}

\begin{defbox}
\textbf{新污染物（Emerging Pollutants）}：由人类活动造成的、目前已在环境中明确存在、危害人体健康和生态环境的，但因其生产使用历史相对较短或危害发现较晚，尚无法律法规和标准予以明确规定或规定不完善的所有在生产建设或其他活动中产生的污染物。一般包括环境激素、抗生素等新化学物质和微/纳塑料等细颗粒物。

\textbf{环境自净（Environmental Self-Purification）}：环境受到污染后，在物理、化学和生物的作用下，对有害物质进行扩散、稀释、氧化还原、生物降解等作用，逐步降低污染物的浓度，减小甚至消除污染物的毒性。

\textbf{公害病（Public Nuisance Disease）}：严重的地区性环境污染性疾病，由政府认定，须经严格鉴定和国家法律的正式认可，有\imp{医学和法律的双重含义}，病人享受国家补贴。

\textbf{毒物兴奋效应（Hormesis Effect）}：某些物质在低剂量时对生物系统具有刺激（有益）作用，而在高剂量时具有抑制（有害）作用，其剂量-效应关系以\imp{双向曲线（倒U或J形）}为特征。典型环境污染物如镉、铅、汞、二噁英等均表现出Hormesis现象。环境因素对人体健康的这种\imp{双重性}（Dual Role）是环境卫生学的重要基础概念。
\end{defbox}

% ----- 1.3 环境问题 -----
\section{全球环境问题与挑战}

\subsection{全球性环境问题}
\begin{enumerate}[leftmargin=*]
    \item \keyword{全球气候变暖}
    \item \keyword{臭氧层破坏}
    \item \keyword{酸雨}
    \item \keyword{生物多样性锐减}
\end{enumerate}

世界第一部关于环境的著作：\imp{《寂静的春天》}（Rachel Carson，1962），标志着现代环保运动的开端。

\textbf{我国环境保护工作的发展历程：}充分理解——没有安全的环境，就不可能有安全的饮水、食品、空气，更不可能有人类健康。

\subsection{世界八大公害事件}

\begin{table}[htbp]
\centering
\caption{世界八大公害事件简表}
\begin{tabularx}{\textwidth}{lXXX}
\hline
\textbf{事件名称} & \textbf{时间} & \textbf{污染物} & \textbf{健康危害} \\
\hline
马斯河谷烟雾事件 & 1930（比利时） & SO$_2$、粉尘 & 呼吸系统疾病 \\
多诺拉烟雾事件 & 1948（美国） & SO$_2$、金属粉尘 & 呼吸系统疾病 \\
伦敦烟雾事件 & 1952（英国） & SO$_2$、烟尘 & 急性中毒、死亡 \\
洛杉矶光化学烟雾 & 1943起（美国） & O$_3$、PAN & 眼睛刺激、呼吸疾病 \\
水俣病 & 1953起（日本） & 甲基汞 & 神经系统损害 \\
痛痛病 & 1955起（日本） & 镉 & 骨骼病变 \\
四日市哮喘 & 1961起（日本） & SO$_2$、粉尘 & 哮喘 \\
米糠油事件 & 1968（日本） & 多氯联苯 & 皮肤病变 \\
\hline
\end{tabularx}
\end{table}

\begin{keybox}
\textbf{水俣病}（Minamata Disease）：由甲基汞污染水体，通过食物链生物放大作用引起慢性甲基汞中毒，主要表现为神经系统损害。

\textbf{痛痛病}（Itai-itai Disease）：由镉污染土壤和水体，通过食物链蓄积引起慢性镉中毒，以骨骼病变和肾损害为特征。
\end{keybox}

% ----- 1.4 复习题 -----
\section{复习题}

\begin{enumerate}[leftmargin=*, itemsep=8pt]
    \item \textbf{【名词解释】}环境卫生学；新污染物（Emerging Pollutants）
    \item \textbf{【名词解释】}一次污染物与二次污染物；环境自净
    \item \textbf{【名词解释】}原生环境与次生环境；公害病
    \item \textbf{【名词解释】}Hormesis效应（毒物兴奋效应）；环境内分泌干扰物（EDCs）
    \item \textbf{【简答题】}简述环境卫生学的研究内容。
    \item \textbf{【简答题】}简述我国环境卫生工作的主要成就及今后任务。
    \item \textbf{【非标准答案题】}试述你对"环境就是民生，青山就是美丽，蓝天也是幸福"的理解，以及环境卫生学在"美丽中国"建设中应发挥的作用。
\end{enumerate}

\subsection*{参考答案要点}

\begin{enumerate}[leftmargin=*, itemsep=5pt]
    \item \textbf{环境卫生学}：研究自然环境和生活环境与人群健康的关系，揭示环境因素对人群健康影响的发生、发展规律，为充分利用环境有益因素和控制环境有害因素提出卫生要求和预防对策，增进人体健康，维护和提高人体健康水平的学科。\textbf{新污染物}：由人类活动造成的、在环境中明确存在但尚无法律法规和标准予以明确规定的污染物，如环境激素、抗生素、微/纳塑料等。

    \item \textbf{一次污染物}：由污染源直接排入环境未发生变化的污染物。\textbf{二次污染物}：一次污染物进入环境后在物理、化学或生物学作用下转化形成的新污染物，其理化性质和毒性通常与原污染物不同。\textbf{环境自净}：环境受污染后，在物理、化学和生物作用下逐步降低污染物浓度、减小或消除污染物毒性的过程。

    \item \textbf{原生环境}：天然形成的未受或少受人为因素影响的环境。\textbf{次生环境}：受人为活动影响形成的环境。\textbf{公害病}：严重的地区性环境污染性疾病，有医学和法律的双重含义。

    \item 环境卫生学的四大研究内容：①环境与健康关系的基础理论研究；②环境因素与健康关系的确认性研究；③创建和引进环境卫生学研究的新技术和新方法；④研究环境卫生监督体系的理论依据。

    \item 我国环境卫生工作的主要成就包括：建立了较为完善的环境卫生标准体系；开展了广泛的环境污染与健康关系的调查研究；积极参与全球环境治理。今后任务：加强基础理论和确认性研究；完善环境卫生法律法规标准体系；加强农村环境卫生工作；开拓新污染物、气候变化与健康等新领域。

    \item （开放性问题，围绕环境保护与人民健康的关系、可持续发展理念、公共卫生工作者的责任等方面展开论述。）
\end{enumerate}

\newpage

% =====================================================================
%  CHAPTER 2: 环境与健康的关系
% =====================================================================
\newpage